\documentclass[11pt]{article}

\usepackage[margin=2.2cm]{geometry}
\usepackage{amsmath,amssymb}
\usepackage{longtable,array,booktabs}
\usepackage{xurl}
\usepackage[colorlinks=true,linkcolor=blue,urlcolor=blue]{hyperref}
\setlength{\emergencystretch}{2em}

\DeclareUrlCommand\leanid{\urlstyle{sf}}
\newcommand{\src}[1]{\path{#1}}
\newcommand{\cD}{c_{\mathrm D}}
\newcommand{\cM}{c_{\mathrm M}}
\newcommand{\Ctpl}{C_{\mathrm{tpl}}}
\newcommand{\Cpad}{C_{\mathrm{pad}}}
\newcommand{\Cbuf}{C_{\mathrm{buf}}}
\newcommand{\Cres}{C_{\mathrm{res}}}
\newcommand{\Csum}{C_{\mathrm{sum}}}
\newcommand{\eps}{\varepsilon}

\title{The Constant in the Two-Dimensional Area Law:\\
    Dependency Chain and Effectivity}
\author{}
\date{2026-10-17}

\begin{document}
\maketitle

\section{The question}

Theorem~1.1 of the two-dimensional area-law manuscript of September~24, 2026
(\src{00-introduction.tex}, lines~39--53, at revision
\texttt{adc7f1241b42e322a6451854ab7e4b4c146bf78a}) states that for every local
dimension $q\ge1$, interaction range $R$, term bound $J>0$ and spectral gap
$\Delta>0$ there is a finite constant $C=C(q,R,J,\Delta)$ such that every gapped
ground state $\Omega$ of an admissible Hamiltonian on a finite induced domain
$\Lambda\subseteq\mathbb Z^2$ and every cut $A\subseteq\Lambda$ satisfy
\[
    S_\Omega(A)\le C\,|\partial_\Lambda A|.
\]
The manuscript asserts only that such a constant exists. It never gives a value:
constants ``may change at different occurrences'' with thresholds ``fixed in the
order stated in the proof'' (\src{01-preliminaries.tex}, lines~134--139), and in
the final assembly ``their magnitude is immaterial'' (\src{10-geometry.tex},
lines~682--683).

This note asks for an explicit function $C_{\rm expl}(q,R,J,\Delta)$ for which
the conclusion of Theorem~1.1 holds, under exactly the hypotheses of
Theorem~1.1. Such a statement is a new auxiliary result. It is stronger than the
source theorem and is not the formalization of any source label, so it belongs
in a blueprint entry of its own.

The note walks the dependency chain of the constant, from the final assembly
down to the four propositions it uses and the collar radius threshold, and
classifies every place where the proof obtains a number without writing it down.
Each such place is one of four kinds.
\begin{itemize}
    \item[(a)] \emph{Eventual thresholds.} A statement ``for all sufficiently
          large $n$'' about an explicit elementary function, from which a
          threshold $N$ is extracted. The manuscript is equally non-quantitative
          at these points (``sufficiently large $n$'': \src{08-scanner.tex},
          lines~198, 215, 412, 739, 817, 849; \src{09-amplification.tex},
          lines~274, 499; \src{10-geometry.tex}, lines~206, 335, 380,
          760--774). Values of convergent series and integrals of explicit
          terms are listed with this kind, since their explicit upper bounds
          use the same tail estimates.
    \item[(b)] \emph{Compactness minima.} A positive constant obtained as the
          minimum of a continuous function on a nonempty compact set.
    \item[(c)] \emph{Replica limits.} A limit as the replica count $k\to\infty$
          used only to pass to the limit in an inequality whose other sides do
          not depend on $k$.
    \item[(d)] \emph{Optimizers and choices.} A point, an index or a minimizer
          that the proof uses as an object. It depends on the domain, the cut or
          the state, and it enters no constant.
\end{itemize}
Only kinds (a) and (b) affect the value of $C$. The verdict, argued in
Section~\ref{sec:verdict}, is that every link of the chain is effective in
principle: each kind (a) threshold is an explicit elementary inequality and each
kind (b) minimum is a positive distance between explicit polygons. No step of
the present proof is ineffective, but the proof does not write the constant
down, and the constant it implicitly defines is very large.

\section{The top level}
\label{sec:top}

The source proof of Theorem~1.1 (\src{10-geometry.tex}, lines~803--857) combines
Proposition~9.6 (the sublinear collar bound), Proposition~10.2 (amplification)
and Proposition~11.2 (the two-family decomposition). The Lean assembly
\leanid{TNLean.PEPS.AreaLaw.uniformAreaLaw_of_collar_amplification_twoFamilies}
(\src{TNLean/PEPS/AreaLaw/BoundaryTheorem.lean}, line~374) fixes the constants in
the following order.

\begin{enumerate}
    \item The template constants are fixed as $D_0=2R+11$ and $\Ctpl=24$
          (the final export \leanid{TNLean.PEPS.AreaLaw.uniformAreaLaw}). Both
          are explicit. The collar proposition requires only $D_0>2R+10$ and
          $\Ctpl\ge24$.
    \item Proposition~9.6 gives a constant $C_1$ and a threshold $N_1$.
    \item Proposition~10.2 is applied with $C_0=\max(C_1,1)$ and decay power
          $M=100$; it gives $C_2$ and $N_2$.
    \item The collar radius threshold
          \leanid{TNLean.PEPS.AreaLaw.Exponents.exists_collar_radius_threshold}
          is applied with coefficient $C_2$ and $\eta=1$; it gives $N_3$.
    \item Proposition~11.2 is applied with lower scale
          $n_*=\max(N_1,N_2,N_3)$; it gives $C_3$.
    \item With $\cD=\max(C_3,0)$ and $\cM=\max(C_2,0)$, the final constant is
          \begin{equation}
              C=\cD\log q+\frac{\cM\,\cD}{2}.
              \label{eq:top}
          \end{equation}
\end{enumerate}
The formula \eqref{eq:top} is explicit. The constant is therefore explicit as
soon as $C_1,N_1,C_2,N_2,N_3$ and $C_3$ are, and the rest of the note concerns
these six numbers.

In Lean, \eqref{eq:top} is \leanid{TNLean.PEPS.AreaLaw.uniformAreaLawConstantOf}
and $n_*$ is \leanid{TNLean.PEPS.AreaLaw.uniformAreaLawScale}; the proof of
Theorem~1.1 at given $C_1,N_1,C_2,N_2,C_3$ is
\leanid{TNLean.PEPS.AreaLaw.regionalEntropy_le_of_collar_amplification_twoFamilies_constants}.
The constant at the named constants of the three propositions, with $D_0=2R+11$
and $\Ctpl=24$, is \leanid{TNLean.PEPS.AreaLaw.uniformAreaLawConstant}, and
\leanid{TNLean.PEPS.AreaLaw.exists_regionalEntropy_le_uniformAreaLawConstant}
proves the bound with it. Its inputs that are not yet closed forms on the
path are listed in Section~\ref{sec:open}: the safe-box constant $C_{\rm box}$
with its entropy bound, which \leanid{TNLean.PEPS.AreaLaw.uniformAreaLaw} takes
from the existential \leanid{TNLean.PEPS.AreaLaw.Scan.exists_boxEntropy_le_boxError};
the truncation constant $C_{\rm tr}$, whose closed form
\leanid{TNLean.PEPS.AreaLaw.latticeTruncationRadiusConstant} is not yet
substituted, and the five transport constants of the scanner, which have no
closed form yet; and the filter-locality constants $(C,c)$ of the positive
constraints, which \leanid{TNLean.PEPS.AreaLaw.uniformAreaLaw} still chooses
existentially although their closed forms are
\leanid{TNLean.PEPS.AreaLaw.latticePositiveConstraintConst} and
\leanid{TNLean.PEPS.AreaLaw.latticePositiveConstraintRate}. The birth-separation constant
is $c_1=1/112$
(\leanid{TNLean.PEPS.AreaLaw.Geometry.isBirthSeparationConstant_birthTemplateSeparationConstant}).

\paragraph{Dependence on the parameters.}
The parameters $q,R,J,\Delta$ enter through four explicit quantities:
\begin{itemize}
    \item the padding scale $Z$ and padding constant $\Cpad$ of the initial
          safe-box estimate (Section~\ref{sec:box}), which carry the dependence
          on $J/\Delta$ and on $q^{O(R^2)}$;
    \item two normalized gaps, one for each filter degree $p$ used in the
          proof. For a filter kernel $f_{p,\Delta/2}$ of degree $p$ put
          $c_p=\max\{1,J(\|f_{p,\Delta/2}\|_1+1)\}$. The normalization and
          the gap depend on $p$ (\src{03-quasilocal.tex}, lines~216--218).
          The scanner chain of Propositions~9.2--9.6 uses $p=1$ and the gap
          $g_1=\Delta/c_1$. The amplification step, Proposition~10.2, uses
          $p_{\rm a}=10^6$, so that $\alpha=p_{\rm a}/(p_{\rm a}+1)>1-10^{-6}$
          (\src{09-amplification.tex}, line~17), and the gap
          $g_{\rm a}=\Delta/c_{p_{\rm a}}$. The two kernels are different
          functions, so $g_1$ and $g_{\rm a}$ are separate inputs of $C$;
    \item the anchor multiplicity $\mu=2^{2R(R+1)}$ and the log-dimension
          $\log q$;
    \item the quasi-local filter constants of Proposition~10.2.
\end{itemize}
The template constant $D_0$ enters only Proposition~11.2, through the layer
width $\lceil4D_0\rceil+2$, which equals $8R+46$ at $D_0=2R+11$, and through
thresholds that grow logarithmically in $D_0$. The collar estimate itself uses
the integer clearance $2R+10$ and does not depend on $D_0$ beyond $D_0>2R+10$.
The template constant $\Ctpl=24$ enters only through the template coefficients
$C_P=\max(18\Ctpl,1)=432$ and $C_B=32\Ctpl=768$ of Proposition~11.2. These
are the named definitions
\leanid{TNLean.PEPS.AreaLaw.Geometry.primaryScannerConstant} and
\leanid{TNLean.PEPS.AreaLaw.Geometry.beltScannerConstant}.

The exponents are fixed rationals:
\[
    e_*=2\cdot10^{-6},\quad
    \eps=\ell_{\rm sc}=10^{-5},\quad
    \mu_{\rm sc}=\tfrac14,\quad
    \nu_{\rm sc}=10^{-3},\quad
    \omega=\tfrac1{20000},\quad
    \kappa_{\rm sc}=10^{-6},
\]
\[
    \beta=1-e_*,\quad
    p=10^6,\quad
    \alpha=\tfrac{p}{p+1},\quad
    \delta_0=2\cdot10^{-7},\quad
    \zeta=1-\tfrac{\delta_0}2,\quad
    \rho_0=\tfrac{1+\beta}2=1-10^{-6},\quad
    \ell_{\rm geo}=10^{-5}.
\]
The bootstrap of Proposition~9.5 uses exponents that depend on the initial
exponent $e_0$ of Section~\ref{sec:box}: $g_0=(1-e_0)/2$, $\mu=1-g_0$,
$\nu=g_0/1000$ and $\ell=\nu/100=(1-e_0)/200000$.

The lower scale $n_*$ enters $C_3$ through the lower layer index $k_0$, which
must satisfy $C_P2^{k_0}\ge n_*$. Since $C_3$ contains the factor $4^{k_0}$
(Section~\ref{sec:geometry}), $\cD$ grows at least like $(n_*/C_P)^2$. The
geometric constant therefore inherits the analytic thresholds $N_1,N_2,N_3$
polynomially.

\section{The dependency tree}
\label{sec:tree}

\subsection{Proposition~9.5: the safe-box estimate}
\label{sec:box}

Proposition~9.5 (\src{08-scanner.tex}, lines~681--805) bounds the safe-box
entropy by $C_{\rm box}\,r^{1+e_*}$ at every scale $r\ge1$
(\leanid{TNLean.PEPS.AreaLaw.Scan.exists_boxEntropy_le_boxError},
\src{TNLean/PEPS/AreaLaw/Scan/PhysicalBoxBootstrap.lean}, line~293). It has two
stages.

\paragraph{Initial estimate (Proposition~3.3).}
The initial estimate (\src{02-initial.tex}, lines~596--669;
\leanid{TNLean.PEPS.AreaLaw.boxEntropy_le_initialBoxConstant_mul_rpow},
\src{TNLean/PEPS/AreaLaw/InitialBoxEstimate.lean}, line~542) is explicit. With
the explicit contour constants $B_{\rm cross}(R)$ and
$K_{\rm budget}(q,R)=1+B_{\rm cross}(R)\log\bigl(e\,q^{1+2R(R+1)}\bigr)^2$
(\leanid{TNLean.PEPS.AreaLaw.bufferCrossConst},
\leanid{TNLean.PEPS.AreaLaw.bufferBudgetConst}),
\begin{align*}
    Z      & =\max\Bigl\{4,\;32B_{\rm cross}(R)\,q^{2(1+2R(R+1))}\frac{J}{\Delta},\;
    128\sqrt{(1+J/\Delta)K_{\rm budget}(q,R)}\Bigr\},                                  \\
    \Cpad  & =\bigl\lceil (R+3)\,e^{(R+1)Z}\bigr\rceil+R+3,\qquad
    p=\frac1{2(2\Cpad+1)^2},                                                            \\
    M      & =8(\Cpad+1),\qquad \lambda=(1-p/2)M^2,                                     \\
    e_0    & =\frac{\log\lambda}{\log M}-1\in(0,1),\qquad
    C_{\rm init}=\Bigl(\log q+\frac{\Cbuf M^2}{\lambda-M}\Bigr)M^{1+e_0},
\end{align*}
where $\Cbuf$ is an explicit expression in the same data and $C_{\rm init}$
is the constant of the scale recurrence
$F(Mr)\le\lambda F(r)+\Cbuf M^2r$ (the QICLean lemma
\leanid{Entropy.exists_le_mul_rpow_of_scale_recurrence}, with $F(1)\le\log q$).
Note that $1-e_0=-\log(1-p/2)/\log M\approx p/(2\log M)$ is of order
$\Cpad^{-2}/\log\Cpad$. These quantities are named:
\leanid{TNLean.PEPS.AreaLaw.initialPaddingScale} ($Z$),
\leanid{TNLean.PEPS.AreaLaw.initialPaddingConst} ($\Cpad$),
\leanid{TNLean.PEPS.AreaLaw.initialBufferConst} ($\Cbuf$),
\leanid{TNLean.PEPS.AreaLaw.initialBoxProb} ($p$),
\leanid{TNLean.PEPS.AreaLaw.initialBoxScale} ($M$),
\leanid{TNLean.PEPS.AreaLaw.initialBoxRatio} ($\lambda$),
\leanid{TNLean.PEPS.AreaLaw.initialBoxExponent} ($e_0$) and
\leanid{TNLean.PEPS.AreaLaw.initialBoxConstantExplicit} ($C_{\rm init}$), the
value of \leanid{TNLean.PEPS.AreaLaw.initialBoxConstant} at the QICLean constant
\leanid{Entropy.scaleRecurrenceConst}. The bootstrap on the path of
Theorem~1.1 does not use this named form yet: it starts from the existential
form \leanid{TNLean.PEPS.AreaLaw.exists_boxEntropy_le_rpow} (issue \#9677).

The values of $M$, $\lambda$ and $p$ above are those of the Lean proof, not of
the source. The source (\src{02-initial.tex}, lines~614--656) orders the
children at random, takes $p_0=(2\Cpad+5)^{-2}$, $M\ge16(\Cpad+2)$ and
$\lambda=(1-p_0/4)M^2$. The Lean proof averages over cyclic rotations of the
colours instead, which gives $p=1/(2(2\Cpad+1)^2)$, $M=8(\Cpad+1)$ and
$\lambda=(1-p/2)M^2$. The statement proved is the same; only these constants
differ.

\paragraph{Bootstrap rounds.}
The exponent is improved from $e$ to $(1-\ell)e$ in each round
(\src{08-scanner.tex}, lines~693--803;
\leanid{TNLean.PEPS.AreaLaw.Scan.exists_uniform_boxEntropy_improvement}). If the
current bound is $C_j r^{1+e}$, the next constant is
\begin{equation}
    C_{j+1}=576\,C_{\rm out}(C_j)+N(C_j)^2\log q+1+\log q,
    \label{eq:round}
\end{equation}
where $C_{\rm out}(C_j)$ is the output coefficient of the rough scanner
(\leanid{TNLean.PEPS.AreaLaw.Scan.CollarScan.exists_sourceBootstrap_entropy_improvement};
named \leanid{TNLean.PEPS.AreaLaw.Scan.bootstrapRestoredBudgetCoefficient} at the
input \leanid{TNLean.PEPS.AreaLaw.Scan.weakRectangleSourceEntropyCoefficient})
and $N(C_j)=\max(N_s,N_g)$ is the largest of its thresholds
(\leanid{TNLean.PEPS.AreaLaw.Scan.CollarScan.sourceBootstrapEntropyThreshold}) and the
rectangle geometry threshold
(\leanid{TNLean.PEPS.AreaLaw.Scan.roundedClearanceThreshold}). One round is
\leanid{TNLean.PEPS.AreaLaw.Scan.boxImprovementConstant}, with threshold
\leanid{TNLean.PEPS.AreaLaw.Scan.boxImprovementThreshold}. The
recursion is the named sequence
\leanid{TNLean.PEPS.AreaLaw.Scan.boxBootstrapConstant}, and the round count is
\leanid{TNLean.PEPS.AreaLaw.Scan.boxBootstrapRounds}, the least $m$ with
$(1-\ell)^m e_0\le e_*$. The last term $\log q$ appears because the rounds
are iterated with threshold $1$: extending the improved bound from scale $1$ on costs the volume
bound once more. The term $N^2\log q$ absorbs the bounded scales by the
volume bound $F(r)\le r^2\log q$
(\src{TNLean/PEPS/AreaLaw/Scan/ExponentBootstrap.lean}, line~39). The threshold
$N(C_j)$ is a maximum of kind (a) thresholds, rows~1--17 and~32--33 of
Table~\ref{tab:sites}, and depends on $C_j$ through the entropy input of the
scanner. The number of rounds is the explicit integer
\[
    m=\Bigl\lceil\frac{\log(e_0/e_*)}{-\log(1-\ell)}\Bigr\rceil,
\]
which is row~32. The constant $C_{\rm box}=C_m$ is therefore a finite recursion
of explicit operations and kind (a) thresholds. The Lean iteration
\leanid{TNLean.PEPS.AreaLaw.Scan.uniform_power_bound_of_boxBootstrapRounds} runs
exactly these $m$ rounds, and $C_{\rm box}$ is
\leanid{TNLean.PEPS.AreaLaw.Scan.physicalBoxBootstrapConstant}
(\leanid{TNLean.PEPS.AreaLaw.Scan.exists_boxEntropy_le_boxError_of_initial}).
Its arguments include the truncation constant $C_{\rm tr}$ and the transport
constants of the scanner, which are not yet closed forms on the path
(Section~\ref{sec:open}), so $C_{\rm box}$ is not yet explicit in Lean.

\subsection{Proposition~9.6: the sublinear collar bound}

Proposition~9.6 (\src{08-scanner.tex}, lines~808--827 and 942--943;
\leanid{TNLean.PEPS.AreaLaw.sublinearCollarBound_real},
\src{TNLean/PEPS/AreaLaw/PhysicalCollarInformation.lean}, line~201) applies the
scanner at the final exponents, with entropy input coefficient
$C_w=4C_{\rm box}+4\log q+1$. The scanner (Proposition~9.2,
\src{08-scanner.tex}, lines~349--533) and the final-scale information estimate
give a constant $C_m$ and a threshold $N_1=\max(N_s,N_g)$
(\leanid{TNLean.PEPS.AreaLaw.Scan.exists_template_final_mutualInformation_of_boxEntropy}),
and $C_1=\max(C_m/2,10)$
(\leanid{TNLean.PEPS.AreaLaw.sublinearCollarBound_of_mutualInformation_le}
replaces the constant $C_m/2$ by $\max(C_m/2,10)$). In Lean these are
\leanid{TNLean.PEPS.AreaLaw.collarBoxConstant} ($C_w$),
\leanid{TNLean.PEPS.AreaLaw.sublinearCollarConstant} ($C_1$) and
\leanid{TNLean.PEPS.AreaLaw.sublinearCollarThreshold} ($N_1$), with $C_m$ and
$N_1$ given by \leanid{TNLean.PEPS.AreaLaw.Scan.collarConstant} and
\leanid{TNLean.PEPS.AreaLaw.Scan.collarThreshold} at $C_{\rm box}$ replaced by
$C_w$. The theorem
\leanid{TNLean.PEPS.AreaLaw.sublinearCollarBound_real_of_boxEntropy} proves
Proposition~9.6 with this $C_1$ and $N_1$ for every safe-box constant
$C_{\rm box}$, and
\leanid{TNLean.PEPS.AreaLaw.sublinearCollarBound_real_of_initial} does so at
$C_{\rm box}$ given by the bootstrap. Both still state the truncation constant
$C_{\rm tr}$ and the transport constants existentially, and $N_1$ depends on
them (Section~\ref{sec:open}).

The output coefficients of the scanner are explicit expressions in the input
coefficient, as recorded in
\src{docs/paper-gaps/arealaw2d_scanner_inputs.tex}: for instance
$C_*=\max\{C_e^{\rm in},1\}$, $C_{\rm rough}=6C_*+13+6\log q$,
$C_{\delta,0}=4(C_1+1)(C_{\rm rough}+2C_t)C_s$, and the final sharp coefficient
$20C_e+62+24\log q+16\sqrt{C_{\rm tr}}(1+C_{\rm dim})$
(\leanid{TNLean.PEPS.AreaLaw.Scan.exists_finalSharp_error_le}; named
\leanid{TNLean.PEPS.AreaLaw.Scan.finalSharpErrorConstant}). The constants
packaged as existentials along this branch are explicit in substance: the
scale-fact constant is $32$, the transport cap is
$1+(2C_r+3)^2\log q$, and the marginal constants are $C_0^{\rm m}=512e(2/g_1)$ (distinct from the
input $C_0=\max(C_1,1)$ of Proposition~10.2) and
$c_0=[32\sqrt{1+2/g_1}]^{-1}$. The thresholds are rows~1--31 of
Table~\ref{tab:sites}. The relation between the safe-box estimate and the
collar bound is recorded in \src{docs/paper-gaps/arealaw2d_safe_box_collar_scope.tex}.

\subsection{Proposition~10.2: amplification}

Proposition~10.2 (\src{09-amplification.tex}, lines~261--576;
\leanid{TNLean.PEPS.AreaLaw.amplificationBound},
\src{TNLean/PEPS/AreaLaw/AmplificationProof.lean}, line~670) is explicit
formulas packaged as existentials, except for one threshold. With $K=4(M+2)$
and the filter constants $c$, $v$ and $C_g$ of the lattice constraint words,
\begin{align*}
    C'  & =\frac{g_{\rm a}/2+2v+12}{c},\qquad
    B=4+\frac{4\,m!}{(c/2)^m},\quad m=\lceil 2/\alpha\rceil,                       \\
    B_r & =\frac6{g_{\rm a}}\Bigl(C_0+1+\frac{2K}{1-\eps}\Bigr)+\frac{K}{1-\eps},\qquad
    C_{\rm rad}=(C'B_r)^{1/\alpha}+1,                                               \\
    C_2 & =\max\bigl\{C_{\rm rad},\,24(1+C_0\log q)\bigr\},\qquad
    N_2=\max\bigl\{2,\lceil C_gC_0B\rceil,N_t\bigr\}.
\end{align*}
These are \leanid{TNLean.PEPS.AreaLaw.amplificationConstant} and
\leanid{TNLean.PEPS.AreaLaw.amplificationThreshold}, built from
\leanid{TNLean.PEPS.AreaLaw.amplificationGap} ($g_{\rm a}$),
\leanid{TNLean.PEPS.AreaLaw.amplificationDecayRate} ($c$),
\leanid{TNLean.PEPS.AreaLaw.amplificationPrefactor} ($C_g$),
\leanid{TNLean.PEPS.AreaLaw.amplificationRadiusConstant} ($C'$),
\leanid{TNLean.PEPS.AreaLaw.amplificationRadiusCoefficient} ($C_{\rm rad}$),
\leanid{TNLean.PEPS.AreaLaw.polyStretchedExpBound} ($B$),
\leanid{TNLean.PEPS.AreaLaw.goodRealizationThreshold} ($\max\{2,\lceil C_gC_0B\rceil\}$)
and \leanid{TNLean.PEPS.AreaLaw.typicalProjectionThreshold} ($N_t$). The
cut-clearance clause of
Proposition~10.2 needs in addition $n\ge N_3$, the collar radius threshold at
$C_2$, so the full conclusion holds for $n\ge\max(N_2,N_3)$. This costs
nothing in Theorem~1.1, whose lower scale $n_*$ already contains $N_3$.
The quasi-local filter constants are named definitions in QICLean
(\leanid{kernelDecayConst}$(\delta)=2\delta e/\pi$ and \leanid{kernelDecayRate}).
The constants $C_2$ and $N_2$ take the tail constants $(C,c)$ of the positive
constraints as arguments. These have the closed forms
\leanid{TNLean.PEPS.AreaLaw.latticePositiveConstraintConst} and
\leanid{TNLean.PEPS.AreaLaw.latticePositiveConstraintRate}, built from the QICLean
constants \leanid{SpectralFilter.positiveConstraintLocalityConst} and
\leanid{SpectralFilter.positiveConstraintLocalityRate}.
\leanid{TNLean.PEPS.AreaLaw.amplificationBound_explicit} proves
Proposition~10.2 at them, with
\leanid{TNLean.PEPS.AreaLaw.amplificationConstantExplicit} and
\leanid{TNLean.PEPS.AreaLaw.amplificationThresholdExplicit}, and
\leanid{TNLean.PEPS.AreaLaw.amplificationBound} is its corollary. The proof of
Theorem~1.1 in \leanid{TNLean.PEPS.AreaLaw.uniformAreaLaw} still takes $(C,c)$
from the existential QICLean statement
\leanid{SpectralFilter.exists_norm_positiveConstraint_sub_map_le}
(Section~\ref{sec:open}).
The typical-projector threshold $N_t$ (row~34) is kind (a): it compares a
stretched exponential $Ce^{-cn^{1/10}}$ with $n^{-2K}$. The constants $c$, $v$,
$C_g$ and $c_{p_{\rm a}}$ involve the $L^1$ norm of the degree-$p_{\rm a}$
filter kernel and a few series of
the form $\sum_n(1+n)^2e^{-cn^\alpha}$ (row~43), which are definite numbers with
explicit upper bounds.
The omission constants $A_o$, $v_o$ and $\gamma$ of the lattice constraint words are
the closed forms \leanid{latticeOmissionPrefactor}, \leanid{latticeOmissionRate} and
\leanid{latticeOmissionExponent} of the localization constants $(C,c)$ of the positive
constraints; they are built from these series and from the row series
$\sum_l(1+l)^4e^{-(b/2)l^\alpha}$ of the event kernel.

\subsection{The collar radius threshold}

The threshold $N_3$ (\src{TNLean/PEPS/AreaLaw/Amplification/RadiusScales.lean},
line~75; row~35) requires
$\lfloor n^{1-\eps}\rfloor+2r<n^\beta$ for every $r\le C_2n^{(1-\eps)/\alpha}$.
The exponent gaps are fixed rationals,
\[
    \beta-(1-\eps)=8\cdot10^{-6},\qquad
    \beta-\frac{1-\eps}{\alpha}=7.00001\cdot10^{-6},
\]
and the threshold is the closed form
\leanid{TNLean.PEPS.AreaLaw.Exponents.collarRadiusThreshold}$(C_2,1)$,
\[
    N_3=\max\Bigl\{2,\,\Bigl\lceil\max\bigl\{1,\,4^{1/(8\cdot10^{-6})},\,
    (8\max(C_2,0))^{1/(7.00001\cdot10^{-6})}\bigr\}\Bigr\rceil\Bigr\}.
\]
Each of the two inner terms has the form $\max\{1,(a/(\eta/4))^{1/d}\}$ with
$\eta=1$, for the coefficients $a=1$ and $a=2\max(C_2,0)$ and the two gaps $d$
above. This is an explicit power bound, $N_3\approx(8C_2)^{1.43\cdot10^5}$.

\subsection{Proposition~11.2: the two-family decomposition}
\label{sec:geometry}

Proposition~11.2 (\src{10-geometry.tex}, lines~98--140, with the constants
fixed in lines~756--774) is assembled in the Lean proof
\leanid{TNLean.PEPS.AreaLaw.isTwoFamilyDecompositionConstant_of_one_le}, and
every constant except the birth-separation constant $c_1$ is a named
definition. The function
\leanid{TNLean.PEPS.AreaLaw.twoFamilyDecompositionConstant}$(D_0,\Ctpl,n_*,c_1)$
is a constant of the decomposition at $n_*$ for every $c_1$ satisfying
\leanid{TNLean.PEPS.AreaLaw.Geometry.IsBirthSeparationConstant}
(\leanid{TNLean.PEPS.AreaLaw.isTwoFamilyDecompositionConstant_twoFamilyDecompositionConstant});
for $\Ctpl<1$ it uses $\max(\Ctpl,1)$. The constants are fixed in the
following order.
\begin{enumerate}
    \item The template coefficients $C_P=\max(18\Ctpl,1)=432$ and
          $C_B=32\Ctpl=768$
          (\leanid{TNLean.PEPS.AreaLaw.Geometry.primaryScannerConstant},
          \leanid{TNLean.PEPS.AreaLaw.Geometry.beltScannerConstant}).
    \item The birth-separation constant $c_1$
          (\leanid{TNLean.PEPS.AreaLaw.Geometry.exists_birthSeparation};
          Lemma~11.5, lines~573--659). It is an absolute constant: it depends on
          none of $q,R,J,\Delta,D_0,\Ctpl$. The source obtains it as a compactness
          minimum (``has a positive minimum \dots\ the list of rescaled
          patterns is finite'', \src{10-geometry.tex}, lines~596--631). It is
          now explicit. For two polygons cut out by belt half-planes, each
          pair of points, one from each polygon, is either at sup distance at
          least $c_2=1/28$, or both points lie in the open quarter-square
          about a common integer contact and satisfy the angular estimate
          (\leanid{TNLean.PEPS.AreaLaw.Geometry.birthSeparationConstant},
          \leanid{TNLean.PEPS.AreaLaw.Geometry.facetPolygon_contact_dichotomy}).
          With the angular constant $c_3=1/2$ and the normalization of belt
          cells to side $4$, this gives $c_1=1/112$
          (\leanid{TNLean.PEPS.AreaLaw.Geometry.birthTemplateSeparationConstant}).
          The assembly of Proposition~11.2 takes $c_1$ as an input
          (\leanid{TNLean.PEPS.AreaLaw.Geometry.IsBirthSeparationConstant}), and
          \leanid{TNLean.PEPS.AreaLaw.Geometry.isBirthSeparationConstant_birthTemplateSeparationConstant}
          shows that $1/112$ is one. The proof of Theorem~1.1 in
          \leanid{TNLean.PEPS.AreaLaw.uniformAreaLaw} uses $c_1=1/112$; the
          stand-alone \leanid{TNLean.PEPS.AreaLaw.twoFamilyDecomposition_of_one_le}
          still takes $c_1$ from the existential
          \leanid{TNLean.PEPS.AreaLaw.Geometry.exists_isBirthSeparationConstant}.
    \item The stopping scale $S_*=\max\{T(1/128),T(1/4),2,U\}$
          (\leanid{TNLean.PEPS.AreaLaw.twoFamilyStoppingScale}), with
          $U=\max\{U_1,U_2,n_*/C_B,1\}$. Here
          $T(a_0)=\max\{1,(1000/a_0)^{1/(1-\rho_0)}\}$
          (\leanid{TNLean.PEPS.AreaLaw.Geometry.repairMarginThreshold}) is the
          removal-margin threshold $d(S)\le(a_0/1000)S$ (row~39). The isolation
          constant $1/128$ of the initial stars enters the templates and the
          separation through the repair induction, and $1/4$ enters the metric
          hypotheses of the counting estimates; $T(1/128)$ is of order
          $128000^{10^6}$, the larger of the two margin thresholds. $U_1$ is the fine-clearance threshold
          $2\lceil C_Bu\rceil^\beta<c_1u^{\rho_0}$. The Lean proof takes
          $U_1=\max\{1,(4(C_B+1)^\beta/c_1)^{1/(\rho_0-\beta)}\}$
          (\leanid{TNLean.PEPS.AreaLaw.Geometry.fineClearanceThreshold}), the
          power-ratio threshold for
          $2(C_B+1)^\beta u^\beta\le(c_1/2)u^{\rho_0}$, with
          $\rho_0-\beta=10^{-6}$. At $C_B=768$ and $c_1=1/112$ its base is
          about $3076/c_1$, so $U_1$ is of order $(3076/c_1)^{10^6}$; $U_2$
          is the belt-replacement threshold
          $\lfloor\lceil C_Bu\rceil^{1-\ell_{\rm geo}}\rfloor\le u$
          (\leanid{TNLean.PEPS.AreaLaw.Geometry.beltReplacementScaleThreshold}),
          and $U$ is \leanid{TNLean.PEPS.AreaLaw.twoFamilyLengthBound}.
    \item The layer width $C=\lceil4D_0\rceil+2$.
    \item The lower layer index
          \leanid{TNLean.PEPS.AreaLaw.twoFamilyLowerLayerIndex}
          $k_0=\max\{5\cdot10^7,K_T,K_U,K_P,K_n\}$, where
          $K_T=\max\{K_\zeta(S_*),K_c,K_P(C_P,1/4)\}$
          (\leanid{TNLean.PEPS.AreaLaw.Geometry.birthTemplateLayerThreshold})
          collects the template, metric and separation thresholds. The metric
          and separation thresholds are both $K_\zeta(S_*)$, the least layer
          whose fine side reaches $S_*$ (row~40), and $K_c$
          (\leanid{TNLean.PEPS.AreaLaw.Geometry.descendantCutClearanceThreshold})
          is the cut-clearance layer of rows~37--38: the dyadic index of the
          power-ratio thresholds for $2R^{\zeta\rho_0}$,
          $8\max(D_0,0)R^\zeta$ and $2(\max(C_B,0)+1)^\beta R^{\zeta\beta}$
          against $R/2$. Each is logarithmic in its data.
    \item The counting constants, named definitions in terms of the series
          value $B(p)=1+\sum_k(k+1)^p2^{-\delta_0k/2}$
          (\leanid{TNLean.PEPS.AreaLaw.Geometry.polynomialDyadicBudget},
          row~41): $M_t=576\,L\,B(L)$
          (\leanid{TNLean.PEPS.AreaLaw.Geometry.repairMarkCountConstant}),
          where $L$ is the descendant-count exponent built from the depth
          constant $D=(|\log\log S_*|+1)/(-\log\rho_0)+1$ through
          $G=\max\{11^{D+1},D\log11\}$, $E=\max\{G2^G,G\}$ and
          $L=\max\{E(\log3+1)^E,E\}$; $M_p=32\,C_P^{-100}B(0)$
          (\leanid{TNLean.PEPS.AreaLaw.Geometry.primaryErrorSumConstant}); and
          $L_q=512\,B(0)+8M_t$
          (\leanid{TNLean.PEPS.AreaLaw.Geometry.repairRetainedCountConstant}).
\end{enumerate}
The constant is then
\begin{align*}
    C_3    & =\max(\Cres,\Csum),                                                    \\
    \Cres & =2\bigl(2(C+1)2^{k_0}+1\bigr)^2+(2S_*+1)^2M_t(2C+1)^2,\qquad
    \Csum=(M_p+L_q)(2C+1)^2.
\end{align*}
It is explicit in $C$, $S_*$ and $k_0$, and $\cD\gtrsim4^{k_0}$. The
descendant-count exponent is very large, roughly $L\approx2^{G2^G}$ with
$G\approx(\log S_*)^{2.4\cdot10^6}$, and $B(L)$ grows like
$L!\,(2/(\delta_0\log2))^{L}$, so $M_t$ is by far the largest factor of
$C_3$. This affects the size of the constant, not its explicitness.

\section{Inventory of extraction sites}
\label{sec:inventory}

\subsection{Elementary shapes}

Every kind (a) threshold in the chain reduces to one of the following
elementary inequalities, each of which has an explicit threshold. Throughout,
$x\ge1$, $A\ge0$, $b>0$ and $u<v$.
\begin{itemize}
    \item[(P1)] \emph{Positive power against a constant}: $b\le n^y$ with
          $y>0$ holds for $n\ge\lceil\max(b,1)^{1/y}\rceil$.
    \item[(P2)] \emph{Power ratio}: $Ax^u\le bx^v$ holds for
          $x\ge(A/b)^{1/(v-u)}$. Finite sums of such terms are handled termwise.
    \item[(P3)] \emph{Power ratio with a logarithmic factor}:
          $Ax^u(\log x)^L\le bx^v$ with $L>0$. Since
          $\log x\le x^\eta/(e\eta)$ for every $\eta>0$, the choice
          $\eta=(v-u)/(2L)$ shows that
          $x\ge\bigl[(A/b)\,(2L/(e(v-u)))^L\bigr]^{2/(v-u)}$ suffices.
    \item[(P4)] \emph{Rounded logarithmic radius against a power}:
          $A\lceil C(\log n)^2\rceil^j\le n^\eps$. For $n\ge3$,
          $\lceil C(\log n)^2\rceil\le(C+1)(\log n)^2$, which reduces this to
          (P3) with $u=0$ and $L=2j$.
    \item[(P5)] \emph{Stretched exponential against a power}:
          $Ce^{-cn^\gamma}\le n^{-s}$. From $e^y\ge y^m/m!$ with
          $m=\lceil(s+1)/\gamma\rceil$, it holds for
          $n\ge(Cm!/c^m)^{1/(\gamma m-s)}$. The same device already appears with
          an explicit constant in
          \leanid{TNLean.PEPS.AreaLaw.exists_one_add_sq_mul_exp_neg_rpow_le}
          and \leanid{TNLean.PEPS.AreaLaw.exists_one_add_pow_mul_exp_neg_rpow_le}.
    \item[(P6)] \emph{Dyadic layer index}: $M(2^k)^a<2^k$ for
          $k>\log_2\max(M,1)/(1-a)$; $M(2^k)^a\le\eps(2^k)^b$ for
          $k\ge\log_2(M/\eps)/(b-a)$; $C_P2^k\ge n_*$ for
          $k\ge\log_2(n_*/C_P)$; and $(k+1)^pe^{-ak}<1$ for
          $k>8p^2/(e^2a^2)$.
    \item[(P7)] \emph{Rounded scales}: $\lfloor x\rfloor\ge x/2$ for $x\ge2$,
          $\lceil x\rceil\le x+1$, and $\lfloor L/(8m)\rfloor\ge B$ once
          $n^{1-\ell-\mu}\ge16B$ and $n^\mu,n^{1-\ell}\ge2$. These reduce to (P1).
    \item[(P8)] \emph{Logarithm and linear}: $\log n\ge1$ for $n\ge3$; and the
          entropy-continuity cost $C\sqrt t(1+Cn^2\log q)\le n^{-247}$ at
          $t\le n^{-500}$, which holds for $n\ge C(1+C\log q)$.
    \item[(P9)] \emph{Archimedean counts}: the number of rounds
          $\lceil\log(e_0/e_*)/(-\log(1-\ell))\rceil$; the least $m$ with
          $2^m>S$, namely $\lfloor\log_2S\rfloor+1$; and
          $\lceil(p+8)/(\kappa/2)\rceil$.
    \item[(P10)] \emph{Series and integral values}:
          $1+\sum_k(k+1)^p2^{-\delta_0k/2}$, $\sum_d(d+1)^{-2}=\pi^2/6$,
          $\sum_n(1+n)^2e^{-cn^\alpha}$, and the filter norm
          $\|f_{1,\Delta/2}\|_1$ and $\|f_{p_{\rm a},\Delta/2}\|_1$. Each has an explicit upper bound: split at
          the (P5) or (P6) threshold, bound the head termwise and the tail by
          a geometric series.
\end{itemize}

\subsection{The eventual statements}

Table~\ref{tab:leaves} lists the statements ``for all sufficiently large $n$''
from which the thresholds are extracted, with their shape. Each is proved by a
limit of an explicit elementary function. Where the threshold is now a named
closed form, the row cites that definition, or the theorem stated from it on.

{\small
\begin{longtable}{>{\raggedright}p{6.2cm}>{\raggedright}p{6.4cm}p{1.6cm}}
\caption{Eventual statements and their elementary shapes. Paths are relative to
    \texttt{TNLean/PEPS/AreaLaw/}.}\label{tab:leaves}                                                                                     \\
\toprule
Declaration (file:line)                                          & Elementary inequality                                                              & Shape \\
\midrule
\endfirsthead
\toprule
Declaration (file:line)                                          & Elementary inequality                                                              & Shape \\
\midrule
\endhead
\leanid{eventually_le_natCast_rpow} (\src{Scan/Budgets.lean}:51) & $b\le n^y$                                                                         & P1    \\
\leanid{eventually_scaleFacts} (\src{Scan/Budgets.lean}:150)
                                                                 & $n\ge2$, $2\le n^\mu$, $2\le n^{1-\ell}$, $48\le n^{1-\ell-\mu}$; constant $32$    & P1, P7 \\
\leanid{eventually_le_K} (\src{Scan/BandCountGrowth.lean}:69)    & $B\le\lfloor\lfloor n^{1-\ell}\rfloor/(8\lfloor n^\mu\rfloor)\rfloor$              & P7    \\
\leanid{eventually_continuityCost_le},
\leanid{eventually_continuityCost_le_of_fixed_coefficient}
(\src{Scan/Budgets.lean}:194, 229)                               & $C\sqrt t(1+Cn^2\log q)\le n^{-247}$ for $t\le An^{-500}$                          & P8    \\
\leanid{eventually_abs_regionalEntropy_sub_le_of_truncation}
(\src{Scan/TargetEntropyTransfer.lean}:86),
\leanid{eventually_regionalEntropy_positiveDepthPrefix_le}
(\src{Scan/PositivePrefixEntropyTransfer.lean}:141),
\leanid{eventually_source_actualTerminal_preEntropy_bounds}
(\src{Scan/SourceActualTerminalEntropy.lean}:95)                 & continuity cost with $C=\max(4,C_{\rm vol})$ and $A=\sqrt{2/g_1}$                   & P8    \\
\leanid{eventually_status_scale_separation},
\leanid{eventually_one_le_source_bandCount}
(\src{Scan/StatusScaleSeparation.lean}:83, 129)                   & $\lceil C_r(\log n)^2\rceil\le n^\kappa$, $8\le n^{\mu-\kappa}$, $8\le n^{1-\ell-\mu}$ & P1, P4 \\
\leanid{eventually_split_counting_scale_separation}
(\src{Scan/BadOldHistoryDecay.lean}:103)                          & the same and $24\le n^{\mu-\kappa}$                                               & P1, P4 \\
\leanid{eventually_const_mul_roundedLogRadius_le_rpow},
\leanid{eventually_roundedLogRadius_le_rpow},
\leanid{eventually_const_mul_roundedLogRadius_cube_le_rpow},
\leanid{eventually_one_le_roundedLogRadius}
(\src{Scan/RadiusAsymptotics.lean}:91, 99, 105, 119)               & $A\lceil C(\log n)^2\rceil^j\le n^\eps$, $j\in\{1,3\}$; $n\ge2$                    & P4    \\
\leanid{eventually_badHistory_bound_le_rpow}
(\src{Scan/BadHistoryAsymptotics.lean}:119),
\leanid{eventually_bad_history_probability_le}
(\src{Scan/PhysicalBadHistoryDecay.lean}:92),
\leanid{eventually_historyProperties}
(\src{Scan/ScanHistories.lean}:343)                              & geometric bad-history tail at most $n^{-p}$, via radius powers and $n\ge8(C_t+1)$ & P4, P9 \\
\leanid{eventually_transportLogDimBound_roundedLogRadius_le}
(\src{Scan/SubsystemDimensionScale.lean}:57)                     & $\log n\ge1$                                                                       & P8    \\
\leanid{ScannerExponents.eventually_mul_a_sqrt_budget_le}
(\src{Scan/SharpMomentScale.lean}:214)                           & $A\,C\sqrt{2C_B}(\log n)^6\le c\,n^\delta$                                         & P3    \\
\leanid{ScannerExponents.principalDensityThreshold},
\leanid{ScannerExponents.eventually_principalDensity_le_rpow}
(\src{Scan/SourceDensityScale.lean}:42, 147)                     & $An^u(\log n)^L\le bn^v$                                                           & P3    \\
\leanid{finalSelectedEnergyThreshold},
\leanid{eventually_final_selectedEnergy_le_rpow}
(\src{Scan/SourceSelectedEnergyPower.lean}:174, 215)             & $16Cn^{-111/125000}(\log n)^{C_l}\le\frac g4n^{-1/2000}$; $n^{-1000}\le\frac g4n^{-1/2000}$ & P3 \\
\leanid{eventually_finalSharp_log_cost_le},
\leanid{eventually_finalSharp_powerEnvelope_le}
(\src{Scan/FinalSharpErrorPower.lean}:59, 124)                    & $An^u(\log n)^L\le n^v$, with $v=62497/62500$                                      & P3    \\
\leanid{ScannerExponents.eventually_bandCount_le_rpow}
(\src{Scan/FinalSharpScalarBounds.lean}:116)                     & $2\le n^\mu$                                                                       & P1, P7 \\
\leanid{rounded_clearance_of_threshold_le}
(\src{Scan/RoundedClearance.lean}:40)                            & $4n^{1-\ell}\le\eta n$ and $20C_r(\log n)^2\le\eta n$, $\eta=1/(4(C_2+1))$, from \leanid{roundedClearanceThreshold} on & P2, P3 \\
\leanid{collar_radius_scale_tendsto_zero}
(\src{Amplification/RadiusScales.lean}:58)                       & $(n^{1-\eps}+2Cn^{(1-\eps)/\alpha})/n^\beta\to0$                                   & P2    \\
\leanid{typicalProjectionThreshold}
(\src{Amplification/TypicalProjectors.lean}:202)                 & $Ce^{-cn^{1/10}}\le n^{-s}$                                                         & P5    \\
\leanid{descendantCutClearanceThreshold},
\leanid{eventually_descendant_cutClearance_scales}
(\src{Geometry/DescendantCutClearance.lean}:95, 179)            & $M(2^k)^a<2^k$ for $a\in\{\zeta\rho_0,\zeta,\zeta\beta\}$                           & P6    \\
\leanid{primaryScaleLimitThreshold},
\leanid{eventually_primaryScaleLimits},
\leanid{le_primaryScannerScale}
(\src{Geometry/PrimaryBirthTemplates.lean}:461, 567, 580)       & $M(2^k)^a\le\eps(2^k)^b$; $C_P2^k\ge n_*$                                          & P6    \\
\leanid{eventually_fine_clearance} (\src{Geometry/RepairDepth.lean}:270)
                                                                 & $2\lceil C_Bu\rceil^\beta<c_1u^{\rho_0}$, gap $\rho_0-\beta=10^{-6}$               & P2    \\
\leanid{eventually_repairHalfSide_le},
\leanid{eventually_repairHalfSide_margins}
(\src{Geometry/SimultaneousRepair.lean}:114, 123)                & $S^{\rho_0}\le(a_0/1000)S$, gap $1-\rho_0=10^{-6}$                                 & P2    \\
\leanid{eventually_beltReplacementScaleLimit}
(\src{Geometry/BeltReplacementTemplates.lean}:326)               & $\lfloor\lceil C_Bu\rceil^{1-\ell}\rfloor\le\eps u$, gap $\ell=10^{-5}$            & P2, P7 \\
\bottomrule
\end{longtable}
}

\subsection{Extraction sites}

Table~\ref{tab:sites} lists every site under \src{TNLean/PEPS/AreaLaw} at which
the proof extracts a number from an eventual statement, an Archimedean bound, a
series, a compactness argument or a least-element choice. The column ``Branch''
records where the number is used: B (the safe-box bootstrap,
Proposition~9.5), S (the scanner at the final exponents, Proposition~9.6),
A (Proposition~10.2), R (the collar radius threshold $N_3$), G
(Proposition~11.2), and ``off'' for a site whose declaration is not used in the
proof of Theorem~1.1. A site used by the bootstrap is also used by the collar
bound, since the latter consumes $C_{\rm box}$. Rows marked $\dagger$ belong to
the Lean proof of Proposition~11.2 and of Theorem~1.1
(\leanid{TNLean.PEPS.AreaLaw.twoFamilyDecomposition_of_one_le},
\leanid{TNLean.PEPS.AreaLaw.uniformAreaLaw}). Line numbers are those of the
main branch on 2026-10-18. Rows~1--31 and~33--35 cite the existential statement
and, after ``named'', the closed-form threshold from which it is now proved; the
existential form is kept as a corollary of the named one. Row~32 cites the round
count itself. Each $\dagger$ row now uses a closed-form
threshold, defined in \src{Geometry/TwoFamilyExplicitThresholds.lean} or, for
the primary scanner scale, in \src{Geometry/PrimaryBirthTemplates.lean}, and
named in its row; \leanid{isTwoFamilyDecompositionConstant_twoFamilyConstant} states
the constant of Proposition~11.2 as \leanid{twoFamilyConstant} of these
thresholds and of the constants of the template, isolation, separation and
counting estimates. Those constants are now named as well (Section~\ref{sec:geometry}),
and rows~37--41 of the geometry block have closed forms; the remaining
input of \leanid{twoFamilyDecompositionConstant}, the birth-separation
constant of rows~57--58, is $c_1=1/112$ on the path of Theorem~1.1.

Sites that merely read a property of an existing witness (for example the inequality $C_1=32\ge1$ read off from the witness
of an eventual statement) are listed
as kind (d); rows~60--63 no longer read anything, since $C_1=32$ is now
passed directly, and are kept so that the row numbers stay fixed. Not listed are the five uses of uniqueness witnesses in
\src{Geometry/ElementarySideReciprocity.lean}, \src{Geometry/InitialSectorRuns.lean},
\src{Geometry/BeltFanColors.lean}, \src{Geometry/RepairLatticeCoverage.lean} and
\src{Geometry/InitialActiveRays.lean}, which select the unique opponent of a
side and extract no number, and the construction of an eventual statement
from a finite graph radius in \src{Amplification/GraphChannelSpatialOmission.lean},
line~59, which extracts nothing.

{\small
\begin{longtable}{r>{\raggedright}p{4.3cm}>{\raggedright}p{6.3cm}cc>{\raggedright\arraybackslash}p{1.3cm}}
\caption{Extraction sites. Paths are relative to
    \texttt{TNLean/PEPS/AreaLaw/}.}\label{tab:sites}                                                                                                               \\
\toprule
\#  & Site and declaration                                      & Elementary statement used                                                    & Kind & Branch & Shape  \\
\midrule
\endfirsthead
\toprule
\#  & Site and declaration                                      & Elementary statement used                                                    & Kind & Branch & Shape  \\
\midrule
\endhead
\multicolumn{6}{l}{\emph{Scanner and bootstrap thresholds}}                                                                                                    \\
1   & \src{Scan/RoughSelectedDensity.lean}:137
\leanid{exists_rough_selectedDensity_principal}             & scale facts and $2\le n^\mu$; named \leanid{ScannerExponents.scaleFactsThreshold}  & a    & B      & P1, P7 \\
2   & \src{Scan/BootstrapPrincipalEnergyBounds.lean}:241
\leanid{exists_bootstrap_principal_energy_bound}            & scale facts, $4C_1\le n^\nu$, $2C_\delta\le n^\nu$, $(\log n)^{C_l}$ against a power of $n$; named \leanid{bootstrapPrincipalEnergyThreshold}  & a & B & P1, P3 \\
3   & \src{Scan/BootstrapReplicaAbsorption.lean}:125
\leanid{exists_bootstrap_comparison_loss_threshold}         & $n^{-\nu/4}<\frac12$ and $C(1+C_e)n^{-\nu/4}<1$; named \leanid{bootstrapComparisonLossThreshold}  & a    & B      & P1     \\
4   & \src{Scan/BootstrapBinaryEntropyCost.lean}:69
\leanid{exists_bootstrap_binEntropy_cost_bound}             & scale facts; named \leanid{ScannerExponents.scaleFactsThreshold}  & a    & B      & P1, P7 \\
5   & \src{Scan/SourceBootstrapEntropy.lean}:448
\leanid{exists_sourceBootstrap_entropy_improvement}         & target entropy continuity cost; named \leanid{CollarScan.sourceBootstrapEntropyThreshold}  & a    & B      & P8     \\
6   & \src{Scan/SourceBootstrapEntropy.lean}:448 (same)          & $4\le n^\nu$                                                                 & a    & B      & P1     \\
7   & \src{Scan/SourceBootstrapEntropy.lean}:448 (same)          & $\log n\ge1$                                                                 & a    & B      & P8     \\
8   & \src{Scan/ScannerRoundedAdmissibility.lean}:58
\leanid{ScannerExponents.exists_rounded_admissibility}      & status scale separation, $2\lceil C_r(\log n)^2\rceil\le n^\kappa$, $K\ge8$; named \leanid{ScannerExponents.roundedAdmissibilityThreshold}  & a    & B      & P1, P4 \\
9   & \src{Scan/FixedWeightTransportSmallness.lean}:85
\leanid{ScannerExponents.exists_fixedWeightFour_transport_smallness} & $n^{1-\ell-\mu}\le CK$, $\log n\ge1$, $4CC_d(\log n)^4\le c_t n^{1-\ell-\mu}$, $K\ge8$; named \leanid{ScannerExponents.fixedWeightTransportThreshold}  & a & B & P3, P7 \\
10  & \src{Scan/RoundedClearance.lean}:104
\leanid{exists_rounded_clearance_threshold}                 & $(4n^{1-\ell}+20(C_r(\log n)^2+1))/n<1/(C_2+1)$; named \leanid{roundedClearanceThreshold}  & a    & B      & P2, P3 \\
11  & \src{Scan/RoundedClearance.lean}:104 (same)                 & $C_2s\ge\max(N,1)$                                                           & a    & B      & P1     \\
12  & \src{Scan/BadOldHistoryDecay.lean}:185
\leanid{exists_badOldHistory_threshold}                     & bad-history tail at most $n^{-P}$; split-counting scales; named \leanid{badOldHistoryThreshold}  & a    & B      & P1, P4 \\
13  & \src{Scan/GenericChargeEnergy.lean}:98
\leanid{exists_generic_badOldHistory_threshold}             & the same at general exponents; named \leanid{ScannerExponents.badOldHistoryThreshold}  & a    & B      & P1, P4 \\
14  & \src{Scan/GenericChargeEnergy.lean}:378
\leanid{ScannerExponents.exists_actualCharge_energy_coefficients} & scale facts and $\log n\ge1$; named \leanid{ScannerExponents.actualChargeEnergyThreshold}  & a    & B      & P7, P8 \\
15  & \src{Scan/ActualChargeEnergyCoefficients.lean}:272
\leanid{exists_actualCharge_energy_coefficients}            & scale facts and $\log n\ge1$; named \leanid{finalActualChargeEnergyThreshold}  & a    & B      & P7, P8 \\
16  & \src{Scan/SourceRoughSelection.lean}:462
\leanid{exists_sourceTarget_selected_actualCharge_rough_comparison} & prefix entropy continuity cost; named \leanid{ScannerExponents.sourceRoughSelectionThreshold}  & a    & B      & P8     \\
17  & \src{Scan/SourceRoughSelection.lean}:462 (same)            & scale facts and $K\ge8$                                                      & a    & B      & P1, P7 \\
18  & \src{Scan/SharpMomentScale.lean}:247
\leanid{ScannerExponents.exists_a_le_moment_radius}         & $(\log n)^6$ against $n^\delta$; named \leanid{ScannerExponents.momentRadiusThreshold}  & a    & S      & P3     \\
19  & \src{Scan/SourceDensityScale.lean}:160
\leanid{ScannerExponents.exists_principalDensity_le_rpow}   & principal density at most $n^{-r}$; named \leanid{ScannerExponents.principalDensityThreshold}  & a    & S      & P3     \\
20  & \src{Scan/SourceEndpointRoughBudget.lean}:129
\leanid{exists_source_actualScan_suppliedPrevector_endpoint_budget} & prefix entropy continuity cost; named \leanid{truncationEntropyThreshold}  & a    & S      & P8     \\
21  & \src{Scan/SourceSelectedSharpEnergy.lean}:393
\leanid{exists_sourceTarget_selected_actualCharge_sharp_energy} & target entropy continuity cost; named \leanid{sourceSelectedSharpEnergyThreshold}  & a    & S      & P8     \\
22  & \src{Scan/SourceSelectedSharpEnergy.lean}:393 (same)       & $4\le n^{\nu_{\rm sc}}$                                                      & a    & S      & P1     \\
23  & \src{Scan/SourceSelectedSharpEnergy.lean}:393 (same)       & $\log n\ge1$                                                                 & a    & S      & P8     \\
24  & \src{Scan/SourceActualInformation.lean}:273
\leanid{exists_sourceTarget_actual_mutualInformation_bound} & terminal entropy continuity cost; named \leanid{sourceSelectedSharpEnergyThreshold}  & a    & S      & P8     \\
25  & \src{Scan/SourceFinalInformation.lean}:227
\leanid{exists_sourceTarget_finalScale_mutualInformation}   & target entropy continuity cost; named \leanid{sourceFinalInformationThreshold}  & a    & S      & P8     \\
26  & \src{Scan/SourceTargetSharpMoments.lean}:240
\leanid{exists_truncated_target_sharp_moments}              & status scale separation; named \leanid{targetSharpMomentThreshold}  & a    & S      & P1, P4 \\
27  & \src{Scan/FinalSharpScalarBounds.lean}:186
\leanid{exists_finalSharp_scale_bounds}                     & $n^{0.74999}\le C_aK\le C_an^{0.74999}$, $2\le n^{1/4000}$; named \leanid{finalSharpScaleThreshold}  & a    & S      & P1, P7 \\
28  & \src{Scan/FinalSharpTransportSmallness.lean}:82
\leanid{exists_finalSharp_transport_smallness}              & $\log n\ge1$; $(C_aC_d/c_t)n^{-37497/50000}(\log n)^4\le1$; named \leanid{finalSharpTransportThreshold}  & a    & S      & P3, P8 \\
29  & \src{Scan/FinalSharpErrorPower.lean}:217
\leanid{exists_finalSharp_error_le}                         & power envelope at most $n^{62497/62500}$; named \leanid{finalSharpErrorThreshold}  & a    & S      & P3     \\
30  & \src{Scan/SourceSelectedEnergyPower.lean}:232
\leanid{exists_final_selectedEnergy_le_rpow}                & selected energy at most $n^{-1/2000}$; named \leanid{finalSelectedEnergyThreshold}  & a    & S      & P3     \\
31  & \src{Scan/FinalScannerAdmissibility.lean}:75
\leanid{exists_finalScanner_admissibility}                  & status scale separation, radius, band count; named \leanid{finalScannerAdmissibilityThreshold}  & a    & S      & P1, P4 \\
\multicolumn{6}{l}{\emph{Archimedean counts}}                                                                                                                   \\
32  & \src{Scan/ExponentBootstrap.lean}:84
\leanid{boxBootstrapRounds}                                 & least $m$ with $(1-\ell)^me_0\le e_*$, a closed form                         & a    & B      & P9     \\
33  & \src{Scan/BadHistoryAsymptotics.lean}:119
\leanid{eventually_badHistory_bound_le_rpow}                & $J\ge(p+8)/(\kappa/2)$; named \leanid{badHistoryThreshold}  & a    & B      & P9     \\
\multicolumn{6}{l}{\emph{Amplification and the collar radius}}                                                                                                  \\
34  & \src{Amplification/TypicalProjectors.lean}:234
\leanid{exists_typicalProjection_failure_le_inv_pow}        & $Ce^{-cn^{1/10}}\le n^{-2K}$; named \leanid{typicalProjectionThreshold}  & a    & A      & P5     \\
35  & \src{Amplification/RadiusScales.lean}:107
\leanid{Exponents.exists_collar_radius_threshold}           & $(n^{1-\eps}+2Cn^{(1-\eps)/\alpha})/n^\beta<\eta$; named \leanid{Exponents.collarRadiusThreshold}  & a    & R      & P2     \\
\multicolumn{6}{l}{\emph{Geometry}}                                                                                                                             \\
36  & \src{Geometry/PolynomialBudget.lean}:61
\leanid{exists_polynomial_dyadic_absorption}                & $(k+1)^pe^{-ak}<1$, $a=\delta_0\log2/4$                                      & a    & G      & P6     \\
37  & \src{Geometry/DescendantCutClearance.lean}:389
\leanid{initialBelt_birthTemplate_cutClearance}             & late-layer scale inequalities; explicit layer \leanid{descendantCutClearanceThreshold} & a & G & P6 \\
38  & \src{Geometry/DescendantCutClearance.lean}:466
\leanid{descendant_birthTemplate_cutClearance}              & late-layer scale inequalities; explicit layer \leanid{descendantCutClearanceThreshold} & a & G & P6 \\
39  & \src{Geometry/RepairInductionMetric.lean}:502
\leanid{exists_repairMetricHypotheses}                      & $S^{\rho_0}\le(a_0/1000)S$; explicit scale \leanid{repairMarginThreshold} & a & G & P2 \\
40  & \src{Geometry/RepairInductionMetric.lean}:502 (same)       & least $m$ with $2^m\ge S_*$; then $K=\lceil m/\zeta\rceil$; explicit layer \leanid{fineScaleIndexThreshold} & a & G & P9 \\
\multicolumn{6}{l}{\emph{Series values}}                                                                                                                        \\
41  & \src{Geometry/PolynomialBudget.lean}:131
\leanid{exists_uniform_polynomial_dyadic_sum_bound}         & $1+\sum_k(k+1)^p2^{-\delta_0k/2}$; named \leanid{polynomialDyadicBudget} & a & G & P10 \\
42  & \src{TruncationSeries.lean}:116
\leanid{sum_shell_le}                                & $\sum_d(d+1)^{-2}=\pi^2/6$; named \leanid{shellSeriesConstant}  & a    & B, S, A & P10   \\
43  & \src{Amplification/GraphOmittedAnchorTail.lean}:353
\leanid{exists_sum_omitted_graphBall_exponentialDistanceProfile_le} & $\sum_n(1+n)^2e^{-cn^\alpha}$ (also in the statements at lines~56 and~227); named \leanid{omittedAnchorTailConstant}  & a & A & P10 \\
\multicolumn{6}{l}{\emph{Thresholds not used in the proof of Theorem~1.1}}                                                                                      \\
44  & \src{Scan/SourceStatusMarginalMoments.lean}:105
\leanid{exists_source_statusMarginal_moment_bounds}         & status scale separation, band count                                          & a    & off    & P1, P4 \\
45  & \src{Scan/SourceChargeSelection.lean}:188
\leanid{exists_sourceTarget_selected_actualCharge_density_energy} & prefix entropy continuity cost                                         & a    & off    & P8     \\
46  & \src{Scan/SourceChargeSelection.lean}:196 (same)           & scale facts and $K\ge8$                                                      & a    & off    & P1, P7 \\
47  & \src{Scan/SourceSelectedSharpDensity.lean}:212
\leanid{exists_sourceTarget_selected_actualCharge_sharp_density} & target entropy continuity cost                                          & a    & off    & P8     \\
48  & \src{Scan/SourceSelectedSharpDensity.lean}:215 (same)      & $4\le n^{\nu_{\rm sc}}$                                                      & a    & off    & P1     \\
49  & \src{Scan/SourceSelectedSharpDensity.lean}:218 (same)      & $\log n\ge1$                                                                 & a    & off    & P8     \\
50  & \src{Scan/SourceTargetSharpUpper.lean}:186
\leanid{exists_truncated_target_actualSharpUpper}           & status scale separation; named \leanid{statusScaleSeparationThreshold}  & a    & off    & P1, P4 \\
51  & \src{Scan/SourceSelectedDensity.lean}:207
\leanid{exists_sourceTarget_selected_actualCharge_source_density} & scale facts                                                            & a    & off    & P1, P7 \\
52  & \src{Scan/SourceSelectedDensity.lean}:209 (same)           & radius at least $1$, $\log n\ge1$, continuity cost                           & a    & off    & P4, P8 \\
53  & \src{Scan/SourceSelectedEnergy.lean}:208
\leanid{exists_sourceTarget_selected_actualCharge_source_density_energy} & $\log n\ge1$                                                    & a    & off    & P8     \\
54  & \src{Scan/SourceScanner.lean}:319
\leanid{scannerStatement}                                   & prefix entropy continuity cost                                               & a    & off    & P8     \\
55  & \src{Scan/BootstrapRoundedClearance.lean}:50
\leanid{exists_two_mul_floor_rpow_ceil_mul_le}              & $2(C_2+1)^\eta\le s^{1-\eta}$                                               & a    & off    & P1     \\
56  & \src{Geometry/ScaleSeparation.lean}:70
\leanid{exists_dyadicScale_side_ratios}                     & least $d$ with $2^d>M$                                                       & a    & off    & P9     \\
\multicolumn{6}{l}{\emph{Compactness minima}}                                                                                                                   \\
57  & \src{Geometry/BirthSeparationAtoms.lean}:567
\leanid{exists_pos_common_contact_dichotomy}                & minimum of $\|x-y\|$ on a compact set of pairs of points of two sectorial polygons away from their common lattice contacts; replaced on the path by $c_2=1/28$ & b & off & -- \\
58  & \src{Geometry/BirthSeparationReplacement.lean}:73
\leanid{exists_pos_le_norm_sub_away_from_intPoints}         & minimum of $\|x-y\|$ on a compact set of pairs away from common integer points; replaced on the path by $c_2=1/28$ & b  & off    & --     \\
\multicolumn{6}{l}{\emph{Replica limits}}                                                                                                                       \\
59  & \src{Scan/ActualRoughLowerRemainder.lean}:135
\leanid{exists_roughLowerRemainder_threshold}               & rough lower remainder below $\eps_{\rm rep}$ for large $k$                   & c    & off    & --     \\
\multicolumn{6}{l}{\emph{Optimizers and choices}}                                                                                                               \\
60  & \src{Scan/SharpMomentScale.lean}:93
\leanid{ScannerExponents.exists_eventually_rpow_le_mul_K}   & $C_1=32$, now written directly; no witness is read                                            & d    & B      & --     \\
61  & \src{Scan/SourceDensityScale.lean}:147
\leanid{ScannerExponents.eventually_principalDensity_le_rpow} & no longer reads $C_1$ from a witness                                          & d    & S      & --     \\
62  & \src{Scan/GenericChargeEnergy.lean}:378
\leanid{ScannerExponents.exists_actualCharge_energy_coefficients} & no longer reads $C_1$ from a witness                                      & d    & B      & --     \\
63  & \src{Scan/ActualChargeEnergyCoefficients.lean}:272
\leanid{exists_actualCharge_energy_coefficients}            & no longer reads $C_1$ from a witness                                            & d    & B      & --     \\
64  & \src{Geometry/InitialMarkFamily.lean}:46
\leanid{finite_support_of_summable_nat_counts}              & a summable count sequence has finite support                                 & d    & G      & --     \\
65  & \src{Geometry/InitialMarkFamily.lean}:169
\leanid{exists_least_incident_scales}                       & least incident belt scale of a mark                                          & d    & G      & --     \\
66  & \src{Geometry/LayerPartition.lean}:78
\leanid{exists_unique_layer_of_not_mem_dyadicNeighborhood}  & least layer containing a point                                               & d    & G      & --     \\
67  & \src{Geometry/CappedDyadicPartition.lean}:132
\leanid{exists_mem_cappedDyadicPartition}                   & largest contained dyadic scale up to the cap                                 & d    & B      & --     \\
68  & \src{Geometry/DyadicExhaustion.lean}:33
\leanid{exists_mem_dyadicNeighborhood}                      & a dyadic neighborhood containing a given point                               & d    & G      & --     \\
69  & \src{Geometry/RepairInductionMetric.lean}:380
\leanid{terminates}                                         & finite termination of each repair branch                                     & d    & G      & --     \\
70  & \src{Geometry/RepairTermination.lean}:45
\leanid{RepairMetricHypotheses.exists_isEmpty_markIndex}    & finite termination of the repair stages                                      & d    & G      & --     \\
71  & \src{Geometry/BirthSeparationAtoms.lean}:72
\leanid{lt_of_mem_interior_intHalfPlane}                    & a point of a neighborhood                                                    & d    & G      & --     \\
72  & \src{Geometry/BirthSeparationAtoms.lean}:274
\leanid{exists_mem_of_finite_cover_ray}                     & a point of a neighborhood                                                    & d    & G      & --     \\
73  & \src{RegularizedPatchMinimum.lean}:132
\leanid{exists_isMinOn_regularizedPatchObjective}           & a minimizing filter on a compact domain (callers at \src{RegularizedPatchLogMinimum.lean}:38, \src{RegularizedPatchMinimum.lean}:240, 317, \src{RegularizedPatchMinimumGrowth.lean}:79) & d & off & -- \\
\multicolumn{6}{l}{\emph{Proposition~11.2 and Theorem~1.1 assembly}}                                                                                            \\
74$\dagger$ & \src{Geometry/RepairActualBirthTemplates.lean}:78
\leanid{repairStage_birth_templates_of_le}         & primary scale limits with $\eps=1/4$; explicit layer \leanid{primaryScaleLimitThreshold} & a & G & P6 \\
75$\dagger$ & \src{Geometry/RepairCharacteristicLengthLowerBound.lean}:55
\leanid{repairStage_positive_rank_length_lower_bound}       & $2^{\lfloor\zeta k\rfloor}\ge U$; explicit layer \leanid{fineScaleIndexThreshold} & a & G & P6 \\
76$\dagger$ & \src{Geometry/TwoFamilyAssembly.lean}:658
\leanid{isTwoFamilyDecompositionConstant_twoFamilyConstant} & fine clearance $U_1$; explicit scale \leanid{fineClearanceThreshold} & a & G & P2 \\
77$\dagger$ & \src{Geometry/TwoFamilyAssembly.lean}:592 (same)           & belt-replacement scale $U_2$; explicit scale \leanid{beltReplacementScaleThreshold} & a & G & P2, P7 \\
78$\dagger$ & \src{Geometry/TwoFamilyAssembly.lean}:486 (same)           & primary scale limits with $\eps=1/8$, giving $K_P$; explicit layer \leanid{primaryScaleLimitThreshold} & a & G & P6 \\
79$\dagger$ & \src{Geometry/TwoFamilyAssembly.lean}:609 (same)           & $C_P2^k\ge n_*$, giving $K_n$; explicit layer \leanid{primaryScannerScaleThreshold} & a & G & P6 \\
\bottomrule
\end{longtable}
}

\paragraph{Summary.}
Table~\ref{tab:sites} has 79 rows: 56 of kind (a), of which 43 are on the path
of Theorem~1.1 and 13 are not; 2 of kind (b); 1 of kind (c), not on the path;
14 of kind (d); and 6 further kind (a) rows in the Lean proof of
Proposition~11.2, now with closed-form thresholds. Rows~37--41 also have
closed forms, and with the named template and counting coefficients the
constant of Proposition~11.2 is explicit: the two kind (b) rows~57--58 are off
the path, replaced by the explicit birth-separation constant
(Section~\ref{sec:geometry}). Every kind (a) row reduces to the
shapes (P1)--(P10). By frequency on the path, the shapes are: power against a
constant and rounded scales (P1, P7), logarithmic factors against power
margins (P3, P4), the entropy-continuity and logarithm bounds (P8), dyadic
layer indices (P6), power ratios with tiny gaps (P2), Archimedean counts (P9),
series values (P10), and a single stretched exponential (P5). The companion
explicit-threshold lemmas are in \src{TNLean/Algebra/ExplicitThresholds.lean}:
\leanid{Real.powerRatioThreshold} (P2),
\leanid{Real.natRpowThreshold} (P1),
\leanid{Real.logPowThreshold} and \leanid{Real.logRpowThreshold} (P3, P4) and
\leanid{Real.stretchedExpThreshold} (P5); the dyadic indices (P6) are
\leanid{TNLean.PEPS.AreaLaw.Geometry.dyadicIndexThreshold} and
\leanid{TNLean.PEPS.AreaLaw.Geometry.fineScaleIndexThreshold}.

\subsection{Replica limits}

The replica count $k$ enters only through limit passages, kind (c). No
threshold in $k$ is extracted on the path of Theorem~1.1 (row~59 is off the
path). The inequality passed to the limit has the form
\[
    (2-C(\tau+E/\tau))\,S\le B_u+B_l+r(k),\qquad r(k)\to0,
\]
and its other sides do not depend on $k$ (\src{08-scanner.tex},
lines~755--769), so the limit gives $S\le B_u+B_l$ without any rate. The
relevant statements are listed below; none contributes to the value of $C$.

{\small
\begin{longtable}{>{\raggedright}p{7.6cm}>{\raggedright\arraybackslash}p{7.4cm}}
\caption{Replica-count limits. Paths are relative to
    \texttt{TNLean/PEPS/AreaLaw/Scan/}.}\label{tab:replica} \\
\toprule
Declaration (file:line)                                                              & Limit used                                                    \\
\midrule
\endfirsthead
\toprule
Declaration (file:line)                                                              & Limit used                                                    \\
\midrule
\endhead
\leanid{entropy_le_of_replica_comparisons} (\src{BootstrapReplicaAbsorption.lean}:42) & passage to the limit in the bootstrap comparison              \\
\leanid{restored_entropy_le_of_selected_rough_comparisons}
(\src{SourceBootstrapEntropy.lean}:47)                                               & the same for the restored entropy                             \\
\leanid{regionalEntropy_le_of_rough_typical_comparisons}
(\src{TypicalReplicaAbsorption.lean}:47)                                             & the same for the typical comparison                           \\
\leanid{information_le_of_eventually_sharpInformationBound}
(\src{SharpInformationScalar.lean}:68)                                               & the same for the sharp information bound                      \\
\leanid{tendsto_log_div_nat} (\src{Selection.lean}:64)                               & $\log(k+2)/k\to0$                                             \\
\leanid{tendsto_roughEndpointRemainder} (\src{ActualEndpointRoughBudget.lean}:40)    & Bell overlap remainder and $C\log(k+2)/k\to0$                 \\
\leanid{tendsto_roughLowerRemainderAtWeight_of_label_dimension},
\leanid{tendsto_roughLowerRemainder_of_label_dimension}
(\src{ActualRoughLowerRemainder.lean}:62, 116)                                       & label-dimension defect $o(k)$ divided by $k$                  \\
\leanid{tendsto_sharpLowerRemainder_of_label_dimension}
(\src{ActualSharpLowerRemainder.lean}:71)                                            & the same for the sharp remainder                              \\
\leanid{tendsto_selectedEnergy_remainder} (\src{SelectedEnergyDelta.lean}:59)        & $(A|\rho_k|^{1/8}+{\rm rem}_k)/{\rm gap}\to0$                 \\
\bottomrule
\end{longtable}
}

\subsection{Dependencies outside this directory}

Five ingredients on the path live in QICLean, at the revision
\texttt{6e5aa9b5} pinned in \src{lake-manifest.json}. The scale-recurrence
constant $(|F(1)|+K/(\lambda-M))M^{\log\lambda/\log M}$ is the named definition
\leanid{Entropy.scaleRecurrenceConst}
(\src{QICLean/Analysis/ScaleRecurrence.lean}, line~69). The quasi-local
filter constants are named definitions, and so are the locality constants
\leanid{SpectralFilter.positiveConstraintLocalityConst} and
\leanid{SpectralFilter.positiveConstraintLocalityRate} of the positive constraints
(\src{QICLean/Analysis/SpectralFilter/FilterLocality.lean}, lines~366 and~372);
the proof of Theorem~1.1 still uses the existential form
\leanid{SpectralFilter.exists_norm_positiveConstraint_sub_map_le} (line~459).
The replica transport estimate
\leanid{TensorPower.ReplicaTransport.transport_finite_family}
(\src{QICLean/Representation/ReplicaTransport/FiniteFamily.lean}, line~39)
supplies the transport constants of the scanner existentially. The typical-tail
threshold
\leanid{Entropy.exists_typical_tail_bound_of_log_pow_twelve_budget}
(\src{QICLean/Analysis/TypicalTailScales.lean}, line~169) extracts a
threshold from a stretched exponential against $n^{-100}$ and from
$(\log n)^6$ against a power, shapes (P5) and (P3). The inventory above is
restricted to \src{TNLean/PEPS/AreaLaw}; a complete explicit constant also
needs the QICLean thresholds on the path, which have the same shapes.

\section{Size of the constant (heuristic)}

\emph{This section is a heuristic from the shape of the recursion. It is not a
line-by-line bound and none of its estimates is proved.}

Write $Z=q^{O(R^2)}J/\Delta$ for the padding scale of Section~\ref{sec:box}.
Then $\Cpad\approx e^{(R+1)Z}$, $1-e_0$ is of order $\Cpad^{-2}$ up to a
logarithm, and $\ell=(1-e_0)/200000$. The number of bootstrap rounds is
$m\approx\log(e_0/e_*)/\ell$, of order $10^6\Cpad^2$, that is $e^{2(R+1)Z}$ up to
a polynomial factor.

In each round the thresholds are of the form (P3) with power gaps of order
$\nu=(1-e_0)/2000$ and coefficients linear in the current constant $C_j$, so
$N(C_j)\approx C_j^{O(1/\nu)}$ and, by \eqref{eq:round},
$C_{j+1}\approx C_j^{O(1/\nu)}$. After $m$ rounds,
\[
    \log C_{\rm box}\approx(\log C_{\rm init})\cdot\nu^{-O(m)},
\]
which is doubly exponential in $m$. With $m\approx e^{2(R+1)Z}$, the box
constant is roughly a tower of three exponentials in $Z$.

The remaining steps add no further level. The collar thresholds are powers of
$C_{\rm box}$ with exponents of order $10^5$; $N_3\approx(8C_2)^{1.43\cdot10^5}$;
and $\cD\gtrsim4^{k_0}\ge(n_*/C_P)^2$, together with absolute factors of order
$(3076/c_1)^{2\cdot10^6/\zeta}$ from the fine clearance $U_1$ and
$128000^{2\cdot10^6/\zeta}$ from the dominant removal margin
$T(1/128)\approx128000^{10^6}$, both entering through $S_*$ and $k_0$. Altogether $C$ is roughly a tower of a few
exponentials in $q^{O(R^2)}J/\Delta$. The tower comes from the bootstrap: its
round count is exponential in $Z$ and each round raises the constant to a large
power.

\section{Verdict and work plan}
\label{sec:verdict}

\paragraph{Verdict.}
The constant of Theorem~1.1 is effective in principle. No step of the present
proof is ineffective. Every kind (a) threshold is an explicit elementary
inequality of one of the shapes (P1)--(P10); the kind (b) minima are positive
distances between polygons with vertices on the quarter mesh and four slopes,
which the mesh argument of \src{10-geometry.tex}, lines~351--363, bounds below
explicitly; the kind (c) limits need no rate; and the kind (d) choices enter no
constant. The constant is not explicit as written, and it is enormous.

\paragraph{Work plan.}
\begin{enumerate}
    \item[2.] \emph{Explicit-threshold lemmas.} Done: the thresholds for the
          shapes (P1)--(P5) are in \src{TNLean/Algebra/ExplicitThresholds.lean}
          (Section~\ref{sec:inventory}), and the dyadic indices (P6) in
          \src{TNLean/PEPS/AreaLaw/Geometry/TwoFamilyExplicitThresholds.lean}.
    \item[3.] \emph{Named witnesses.} Largely done. The statements of the form
          ``there are $C$ and $N$ with $P(C,N)$'' along the path now have named
          constants and thresholds with a theorem at the named values, through
          the scanner (Propositions~9.2, 9.5, 9.6), amplification
          (Proposition~10.2) and geometry (Proposition~11.2); rows~1--35 and
          74--79 of Table~\ref{tab:sites} name them. The inputs that are not
          yet closed forms on the path are listed below.
    \item[4.] \emph{An explicit birth-separation constant.} Done: the two
          compactness minima (rows~57--58) are replaced by $c_2=1/28$, from an
          error bound for polygons cut out by belt half-planes, which gives
          $c_1=1/112$. The compactness lemmas of rows~57--58 remain in the
          library but no proof on the path of Theorem~1.1 uses them.
    \item[5.] \emph{The explicit theorem.} State Theorem~1.1 with
          $C=C_{\rm expl}(q,R,J,\Delta)$ and give it its own blueprint entry,
          now that Proposition~11.2 and Theorem~1.1 are proved on main. A
          recursively defined computable $C_{\rm expl}$ needs no
          further mathematics. An elementary closed form additionally needs a
          bound for the bootstrap recursion \eqref{eq:round}.
\end{enumerate}
Apart from item~4, now done, only the closed-form bound for the recursion is
new quantitative mathematics; the rest is restatement.

\paragraph{Open dependencies.}
\label{sec:open}
The following inputs of $C$ are not yet closed forms on the path of
Theorem~1.1. Until they are, $C_{\rm box}$, $N_1$, $C_2$ and $N_2$ are named
functions of existentially chosen arguments, not explicit numbers.
\begin{itemize}
    \item The transport constants $(c_t,C_{\rm ent},e_{\rm ent},C_{\rm en},e_{\rm en})$
          of the scanner, which satisfy
          \leanid{TNLean.PEPS.AreaLaw.Scan.CollarScan.FiniteFamilyTransportBound}
          and are obtained from the QICLean transport estimate through
          \leanid{TNLean.PEPS.AreaLaw.Scan.CollarScan.exists_finiteFamilyTransportBound}.
          No closed form is available yet (issue \#9703). They enter
          $C_{\rm box}$ and $N_1$.
    \item The truncation constant $C_{\rm tr}$. Its closed form is
          \leanid{TNLean.PEPS.AreaLaw.latticeTruncationRadiusConstant}, but the
          scanner, bootstrap and collar theorems still state it
          existentially (issue \#9705). It enters $C_{\rm box}$ and $N_1$.
    \item The tail constants $(C,c)$ of the positive constraints in
          Proposition~10.2. Their closed forms are
          \leanid{TNLean.PEPS.AreaLaw.latticePositiveConstraintConst} and
          \leanid{TNLean.PEPS.AreaLaw.latticePositiveConstraintRate}, and
          \leanid{TNLean.PEPS.AreaLaw.amplificationBound_explicit} proves
          Proposition~10.2 at them (pull request \#9711), but
          \leanid{TNLean.PEPS.AreaLaw.uniformAreaLaw} still chooses them
          existentially. They enter $C_2$ and $N_2$.
    \item The initial constant $C_{\rm init}$. Its closed form is
          \leanid{TNLean.PEPS.AreaLaw.initialBoxConstantExplicit}, but the
          bootstrap on the path still starts from the existential form
          (issue \#9677). It enters $C_{\rm box}$.
\end{itemize}
Of these, only the transport constants lack a closed form; the other three
need only the substitution of an existing definition.

\end{document}
